1. 목적함수의 비용항목(수송비, 진입비 등)을 어떠한 근거로 정하는가?
-> A.cost_selection_rationale에 정리

※ $c^+$ 관련. 비용 중 패널티 비용은 실제로 현실에서 비용으로 계산하기 어려운 것들을 대표 숫자 하나로 통틀어서 정하는 것

2. baseline model을 검증하기 위한 numerical experiment에서 목적함수 내 기호들은 어떤 자료를 근거로 어떻게 산출되었나?
-> B.parameter_rationale_baseline에 정리

3. ER-DRO 목적함수 생성을 위해 참고한 문헌들은?
-ambiguity set 정하는 방식 참고: CVaR형 방식
($\tau$ 사용 원리 참고)
-radius(반경) 정하는 방식 참고
-수요의 regression 방식 사용:

-------------------------------------
0. 먼저 정할 것: 회귀의 종속변수 $y$가 무엇인가

지금 모델은 $D=M\cdot p$인데 관측되는 것은 통행 수 $D$뿐이라, $M$과 $p$를 따로 회귀하려면 각각의 관측치가 있어야 합니다. 현실적인 분해는:

$M$ (시장 규모): 로보택시 이전부터 있던 기존 ride-hailing 통행량을 관측치로 삼아 도시 특성에 회귀. 데이터가 많음(여러 도시 × 여러 해).
$p$ (로보택시 선택률): 로보택시가 운영 중인 도시에서 로보택시 통행 / 전체 ride-hailing 통행 비율을 관측치로 삼아 회귀. 데이터가 매우 적음(운영 도시 몇 개 × 몇 분기).

이 분해가 안 되면 $D$(또는 인구당 $D$)를 통째로 회귀하는 것이 차선이고, 그 경우 baseline의 $M\cdot p$ 구조를 contextual 버전에서 유지할 필요가 있는지 교수님과 정해야 합니다. 지금 문서(B-1)에는 $M_c$ 고정·$p$ 회귀로 되어 있는데, 그러면 $p$의 관측치가 무엇인지가 첫 질문입니다.

2. 첫 시도로 권할 만한 구성

$M$ 회귀 (인구당 ride-hailing 통행 = 종속변수, 정적 도시 특성):

인구밀도, 무차량 가구 비율, 소득, 일자리 밀도, 공항 이용객, 대중교통 공급(SLD) — 표에 있는 변수 중 ACS·SLD·FAA로 후보 도시 값이 다 구해지는 것들. 인구는 종속변수 정규화에 쓰고 공변량에서는 뺍니다.
시간·요일·날씨는 연 단위 진입 결정에서는 평균에 묻히므로 제외(연도 추세만 하나).

$p$ 회귀 (로보택시 점유율 = 종속변수):

운영 개시 후 경과 기간(경험 축적), 18~44세 비율 또는 대졸 비율(수용성 대리), 기존 ride-hailing 이용 강도(습관), 여기까지 3개 정도.
요금·대기시간·차량 내 기능은 우리가 정하는 것이라 제외하고 ④에서 다룹니다. 신뢰·태도 설문은 후보 도시 값이 없어 제외.

3. 최종 선택은 이론이 아니라 교차검증으로

Kannan 실험대로 소표본에서는 Lasso가 낫습니다. 후보 5~8개를 넣고 Lasso가 남기는 변수를 보고하면 "왜 이 변수냐"에 대한 답이 데이터에서 나옵니다. 그 뒤 두 가지를 반드시 확인합니다.

잔차가 공변량과 무관한가(동분산). 통행 수가 아니라 비율·인구당 값을 종속변수로 두는 이유입니다.
후보 도시의 $x_c$가 학습 데이터 범위 안에 있는가. 범위 밖이면 $\hat f(x_c)$의 예측오차(가정 2)가 커져 반경을 키워야 하는데, 사실상 외삽입니다.

4. 교수님의 관점과 맞추기

9.25 미팅에서 교수님은 발사장 사례로 "어떤 핵심 context가 새 위치를 열게 만드는가"를 말할 수 있어야 한다고 했습니다(적도 vs 극궤도 발사 비율, 공급업체 근접성). 우리에게 대응되는 질문은 "어떤 도시 특성이 진입을 결정하는가"이고, 그 답이 되려면 공변량이 적고 해석 가능해야 합니다. 변수 10개짜리 회귀보다 "무차량 가구 비율과 밀도가 높으면 진입" 같은 3개짜리 정책이 논문의 결론이 됩니다. 이 기준으로 보면 M 쪽에서 무차량 가구·밀도·공항, p 쪽에서 경과 기간·연령 구조 정도가 남을 가능성이 큽니다.

정리하면 순서는: 종속변수 확정 → 표에서 네 기준으로 거르기 → 인구당·비율로 정규화 → Lasso로 최종 선택 → 잔차·외삽 점검. 첫 단계(무엇을 $p$의 관측치로 볼 것인가)가 정해지면 나머지는 기계적으로 진행됩니다.